% ============================================================
%  KOMOREBI CAFÉ — Documentación Oficial TFG
%  Alumna : Irene Reyes Carmona
%  Centro  : I.E.S. Ágora
%  Tutora  : María Lourdes Calleja Flores
%  Curso   : 2025-2026
%  Motor   : XeLaTeX + fontspec · compilar con -shell-escape
% ============================================================
\documentclass[12pt,a4paper,twoside,openright]{report}

% ─── Codificación y tipografía ───────────────────────────────
\usepackage{fontspec}
\setmainfont{Times New Roman}
\setsansfont{Arial}
\setmonofont[Scale=0.9]{Courier New}

% ─── Soporte CJK (caracteres japoneses) ──────────────────────
\usepackage{xeCJK}
\setCJKmainfont{Yu Gothic}
\setCJKsansfont{Yu Gothic}
\setCJKmonofont{Yu Gothic}

% ─── Idioma ──────────────────────────────────────────────────
\usepackage[spanish,es-tabla,es-noshorthands]{babel}
\usepackage{csquotes}

% ─── Márgenes y geometría ────────────────────────────────────
\usepackage[a4paper,
  top=2.5cm, bottom=2.5cm,
  left=3cm,  right=2.5cm,
  headheight=15pt]{geometry}

% ─── Encabezados y pie de página ─────────────────────────────
\usepackage{fancyhdr}
\pagestyle{fancy}
\fancyhf{}
\fancyhead[LE,RO]{\thepage}
\fancyhead[RE]{\textit{Komorebi Café}}
\fancyhead[LO]{\textit{\leftmark}}
\renewcommand{\headrulewidth}{0.4pt}

% ─── Tipografía avanzada ─────────────────────────────────────
\usepackage{microtype}
\usepackage{setspace}
\onehalfspacing

% ─── Colores ─────────────────────────────────────────────────
\usepackage[dvipsnames,table]{xcolor}
\definecolor{komorebi}{HTML}{6B8F71}      % verde café
\definecolor{komorebi-dark}{HTML}{3D5A42}
\definecolor{komorebi-light}{HTML}{EBF3EC}
\definecolor{code-bg}{HTML}{F8F8F8}
\definecolor{link-color}{HTML}{1A5276}

% ─── Títulos de capítulo personalizados ──────────────────────
\usepackage{titlesec}
\titleformat{\chapter}[display]
  {\normalfont\Large\bfseries\color{komorebi-dark}}
  {\chaptertitlename\ \thechapter}{10pt}{\LARGE}
\titlespacing*{\chapter}{0pt}{-20pt}{20pt}

% ─── Listas ──────────────────────────────────────────────────
\usepackage{enumitem}

% ─── Tablas ──────────────────────────────────────────────────
\usepackage{booktabs}
\usepackage{longtable}
\usepackage{tabularx}
\usepackage{multirow}
\usepackage{array}
\usepackage{pdflscape}

% ─── Figuras ─────────────────────────────────────────────────
\usepackage{graphicx}
\usepackage{float}
\usepackage{caption}
\usepackage{subcaption}
\usepackage{svg}           % requiere Inkscape en PATH
\svgsetup{inkscapepath=svgdir}

% ─── Código fuente ───────────────────────────────────────────
\usepackage{minted}
\usemintedstyle{friendly}
\setminted{
  fontsize=\small,
  linenos=true,
  breaklines=true
  % bgcolor, frame y framesep se controlan por tcolorbox en los entornos phpcode/sqlcode/bashcode
}
\usepackage{tcolorbox}
\tcbuselibrary{minted,skins,breakable}
\newtcblisting{phpcode}[2][]{
  listing engine=minted,
  minted language=php,
  title=#2,
  colback=code-bg,
  colframe=komorebi,
  fonttitle=\bfseries,
  breakable,
  #1
}
\newtcblisting{sqlcode}[2][]{
  listing engine=minted,
  minted language=sql,
  title=#2,
  colback=code-bg,
  colframe=komorebi-dark,
  fonttitle=\bfseries,
  breakable,
  #1
}
\newtcblisting{bashcode}[2][]{
  listing engine=minted,
  minted language=bash,
  title=#2,
  colback=code-bg,
  colframe=gray,
  fonttitle=\bfseries,
  breakable,
  #1
}

% ─── Diagramas TikZ ──────────────────────────────────────────
\usepackage{tikz}
\usetikzlibrary{shapes,arrows.meta,positioning,calc,fit,backgrounds,shadows}

% ─── Diagramas de secuencia UML ──────────────────────────────
\usepackage{pgf-umlsd}

% ─── Árboles de directorios ──────────────────────────────────
\usepackage{forest}

% ─── Referencias cruzadas ────────────────────────────────────
\usepackage[hidelinks,
  colorlinks=false,
  pdfauthor={Irene Reyes Carmona},
  pdftitle={Komorebi Café — Documentación TFG},
  pdfsubject={Proyecto de Fin de Grado},
  pdfkeywords={PHP,MVC,Docker,PSR-7,Café temático}
]{hyperref}
\usepackage[spanish]{cleveref}

% ─── Glosario ────────────────────────────────────────────────
\usepackage[abbreviations]{glossaries-extra}
\makenoidxglossaries

% ─── Bibliografía ────────────────────────────────────────────
\usepackage[backend=biber,style=ieee,sorting=none]{biblatex}
\addbibresource{referencias.bib}

% ─── Apéndices ───────────────────────────────────────────────
\usepackage[titletoc]{appendix}

% ─── Cajas informativas ──────────────────────────────────────
\newtcolorbox{infobox}[1][]{
  colback=komorebi-light,colframe=komorebi,
  fonttitle=\bfseries,title=Nota,#1
}
\newtcolorbox{warnbox}[1][]{
  colback=yellow!10,colframe=orange!70!black,
  fonttitle=\bfseries,title=Atención,#1
}

% ─── Glosario: términos y acrónimos ──────────────────────────
\newabbreviation{mvc}{MVC}{Model-View-Controller}
\newabbreviation{psr}{PSR}{PHP Standards Recommendation}
\newabbreviation{dto}{DTO}{Data Transfer Object}
\newabbreviation{csrf}{CSRF}{Cross-Site Request Forgery}
\newabbreviation{rbac}{RBAC}{Role-Based Access Control}
\newabbreviation{api}{API}{Application Programming Interface}
\newabbreviation{rest}{REST}{Representational State Transfer}
\newabbreviation{jwt}{JWT}{JSON Web Token}
\newabbreviation{orm}{ORM}{Object-Relational Mapping}
\newabbreviation{ci}{CI/CD}{Continuous Integration / Continuous Delivery}
\newabbreviation{kds}{KDS}{Kitchen Display System}
\newabbreviation{lcov}{LCOV}{Linux Test Coverage}
\newglossaryentry{komorebi}{
  name={komorebi},
  description={Término japonés (木漏れ日) que describe la luz del sol que se filtra
    entre las hojas de los árboles, símbolo del concepto estético del proyecto.}
}
\newglossaryentry{result-pattern}{
  name={Result pattern},
  description={Patrón de diseño donde los métodos devuelven un objeto \texttt{Result}
    con \texttt{ok} (bool), \texttt{data} y \texttt{message}, en lugar de lanzar
    excepciones para fallos esperados.}
}
\newglossaryentry{idempotency}{
  name={idempotencia},
  description={Propiedad de una operación que produce el mismo resultado aunque se
    ejecute múltiples veces. En la API REST de Komorebi, las reservas usan un
    \texttt{Idempotency-Key} para evitar duplicados en caso de reintentos.}
}
\newglossaryentry{12factor}{
  name={12-Factor App},
  description={Metodología de desarrollo de software que define doce principios
    para construir aplicaciones SaaS escalables, portables y resilientes.}
}

% ============================================================
%  DOCUMENTO
% ============================================================
\begin{document}

% ─── PORTADA ─────────────────────────────────────────────────
\begin{titlepage}
  \centering
  \vspace*{1.5cm}

  % Logo del café
  \includesvg[width=5cm]{images/logos/komorebi-logo-stacked}\\[1.5cm]

  {\fontsize{22}{26}\selectfont\bfseries\color{komorebi-dark}
    Komorebi Café\\[0.4cm]
    \fontsize{18}{22}\selectfont\mdseries
    Plataforma de gestión integral\\
    para cadena de cafés temáticos\\[2cm]}

  {\large Documentación de Proyecto Final\\[0.3cm]
   Ciclo Formativo de Grado Superior —
   Desarrollo de Aplicaciones Web}\\[2.5cm]

  \begin{tabular}{rl}
    \textbf{Alumna:}         & Irene Reyes Carmona\\[4pt]
    \textbf{Centro:}         & I.E.S. Ágora\\[4pt]
    \textbf{Tutora:}         & María Lourdes Calleja Flores\\[4pt]
    \textbf{Año académico:}  & 2025-2026\\
  \end{tabular}

  \vfill

  \includesvg[width=2.5cm]{images/logos/komorebi-logo-icon}\\[0.5cm]
  {\small\color{gray} Komorebi Café · Todos los derechos reservados · 2026}
\end{titlepage}

% ─── FRONTMATTER ─────────────────────────────────────────────
\pagenumbering{roman}
\pagestyle{plain}

% ── Resumen en español ───────────────────────────────────────
\chapter*{Resumen}
\addcontentsline{toc}{chapter}{Resumen}

Komorebi Café es una plataforma web de gestión integral diseñada para una cadena ficticia
de catorce cafés temáticos japoneses. El nombre del proyecto proviene del término japonés
\textit{komorebi} (木漏れ日), que describe la luz solar filtrada entre las hojas de los
árboles, evocando la atmósfera serena y cuidada que caracteriza estos espacios.

El sistema cubre la totalidad del ciclo operativo: gestión de reservas con control de
idempotencia, carta digital con filtros por alérgeno, programa de fidelización, panel de
cocina en tiempo real (\gls{kds}), control de salud de animales del local, operaciones de
recepción con check-in/checkout y generación asíncrona de facturas en PDF.

Desde el punto de vista técnico, la solución está construida sobre PHP~8.4 con un
framework \gls{mvc} propio (sin Laravel ni Symfony como capa de aplicación), siguiendo
los estándares \gls{psr}-7, PSR-14 y PSR-15. La arquitectura sigue los principios
12-Factor~App, con despliegue mediante Docker Compose y Railway. El sistema incluye
más de 250~rutas, 53~servicios de negocio, 29~repositorios y más de
95~endpoints de \gls{api} \gls{rest}.

\bigskip
\textbf{Palabras clave:} PHP~8.4, MVC, Docker, PSR-7, PSR-15, Result pattern,
reservas, fidelización, cafés temáticos, 12-Factor.

\vspace{2cm}

% ── Resumen en inglés ────────────────────────────────────────
\chapter*{Abstract}
\addcontentsline{toc}{chapter}{Abstract}

\textit{Komorebi Café} is a comprehensive web management platform designed for a
fictional chain of fourteen Japanese-themed cafés. The project name derives from the
Japanese term \textit{komorebi} (木漏れ日), describing sunlight filtering through tree
leaves, evoking the serene and carefully crafted atmosphere that defines these spaces.

The system covers the complete operational cycle: reservation management with idempotency
control, a digital menu with allergen filters, a loyalty programme, a real-time Kitchen
Display System (KDS), animal health tracking, reception operations with check-in/checkout,
and asynchronous PDF invoice generation.

Technically, the solution is built on PHP~8.4 with a custom MVC framework (no
Laravel/Symfony application layer), following PSR-7, PSR-14 and PSR-15 standards.
The architecture adheres to the 12-Factor App principles, deployed via Docker Compose
and Railway. The system comprises over 250 routes, 53 business services, 29 repositories
and more than 95 REST API endpoints.

\bigskip
\textbf{Keywords:} PHP~8.4, MVC, Docker, PSR-7, PSR-15, Result pattern, reservations,
loyalty, themed cafés, 12-Factor.

% ── Agradecimientos ──────────────────────────────────────────
\chapter*{Agradecimientos}
\addcontentsline{toc}{chapter}{Agradecimientos}

A María Lourdes Calleja Flores, tutora de este proyecto, por su orientación constante,
su rigor técnico y su disponibilidad a lo largo de todo el proceso de desarrollo y
documentación.

A mis compañeros y compañeras del ciclo, por los debates técnicos, las horas compartidas
y el apoyo mutuo que han hecho posible completar este proyecto.

Y a todas las personas que, como la luz entre las hojas, iluminan los rincones del camino
cuando más lo hace falta.

% ── Índices ──────────────────────────────────────────────────
\tableofcontents
\listoffigures
\listoftables

% ─── MAINMATTER ──────────────────────────────────────────────
\clearpage
\pagenumbering{arabic}
\pagestyle{fancy}

% ============================================================
%  CAPÍTULO 1 — INTRODUCCIÓN GENERAL
% ============================================================
\chapter{Introducción General}
\label{chap:intro}

\section{Origen y concepto del proyecto}
\label{sec:origen}

El punto de partida de Komorebi Café es la pregunta: ¿puede una aplicación web de gestión
ser, al mismo tiempo, un reflejo fiel de la identidad de la marca que sirve? El proyecto
nace de la voluntad de demostrar que la tecnología y la estética no son opuestos: un
sistema de software robusto puede encarnar la filosofía tranquila, detallista y contemplativa
que inspira el concepto de \textit{komorebi}.

La palabra japonesa \textit{komorebi} (木漏れ日) se compone de tres kanjis:
\begin{itemize}[nosep]
  \item 木 (\textit{ki}) — árbol
  \item 漏れ (\textit{more}) — filtrar, escapar
  \item 日 (\textit{bi}) — luz, sol
\end{itemize}

Juntos describen ese fenómeno visual fugaz en que la luz solar se cuela entre las hojas
creando patrones de claroscuro. Es una palabra sin traducción exacta al español, pero
precisamente esa intraducibilidad la convierte en símbolo perfecto: algunas experiencias
solo se pueden \textit{vivir}, no describir.

\section{Los catorce cafés temáticos}
\label{sec:cafes}

La cadena ficticia está compuesta por catorce cafés, cada uno inspirado en un universo
estético japonés diferente. La \cref{tab:cafes} recoge los establecimientos registrados
en el sistema.

\begin{longtable}{@{}lll@{}}
  \caption{Catálogo de cafés temáticos de Komorebi Café}
  \label{tab:cafes}\\
  \toprule
  \textbf{Slug} & \textbf{Nombre} & \textbf{Animal / Categoría}\\
  \midrule
  \endfirsthead
  \multicolumn{3}{l}{\footnotesize\textit{(continuación de la \cref{tab:cafes})}}\\
  \toprule
  \textbf{Slug} & \textbf{Nombre} & \textbf{Animal / Categoría}\\
  \midrule
  \endhead
  \bottomrule
  \endfoot
  neko-no-niwa       & Neko no Niwa (猫の庭)         & Gato (lounge)\\
  usagi-paradise     & Usagi Paradise (うさぎパラダイス) & Conejo (lounge)\\
  soft-cloud         & Soft Cloud (ソフトクラウド)     & Chinchilla (lounge)\\
  chipmunk-forest    & Chipmunk Forest (シマリスの森)  & Ardilla (lounge)\\
  mame-shiba-cafe    & Mame Shiba Café (豆柴カフェ)   & Perro Shiba Inu (playroom)\\
  mipig-cafe         & Mipig Cafe (マイピッグカフェ)   & Cerdito miniatura (playroom)\\
  parrot-talk        & Parrot Talk (おしゃべりオウム)  & Loro (playroom)\\
  capyba-land        & Capyba Land (カピバランド)      & Capybara (farm)\\
  alpaca-hill        & Alpaca Hill (アルパカの丘)      & Alpaca (farm)\\
  little-hooves      & Little Hooves (小さなひづめ)   & Caballo miniatura (farm)\\
  quack-club         & Quack Club (クワッククラブ)     & Pato (farm)\\
  pui-pui-house      & Pui Pui House (プイプイハウス)  & Cobaya (zen)\\
  prairie-town       & Prairie Town (プレーリータウン) & Perrito de pradera (zen)\\
  slow-life          & Slow Life (スローライフ)        & Tortuga (zen)\\
\end{longtable}

\section{Misión, visión y valores}
\label{sec:valores}

\begin{description}[style=nextline]
  \item[Misión]
    Ofrecer a cada cliente una experiencia de visita única, combinando una carta
    artesanal con un entorno temático cuidado y un servicio personalizado habilitado
    por tecnología.
  \item[Visión]
    Ser la referencia en España de la hostelería temática de inspiración japonesa,
    reconocidos tanto por la calidad de nuestros productos como por la excelencia
    tecnológica de nuestros sistemas de gestión.
  \item[Valores]
    \textit{Ma} (間, el espacio entre las cosas), \textit{omotenashi} (hospitalidad
    total), sostenibilidad, belleza en los detalles y apertura cultural.
\end{description}

\section{Alcance y finalidad del proyecto}
\label{sec:alcance}

El proyecto abarca el diseño, implementación, pruebas y documentación de una plataforma
web completa que gestiona:

\begin{enumerate}[nosep]
  \item La presencia pública de la cadena (carta, reservas en línea, historia, FAQ)
  \item La autenticación y gestión de cuentas de usuario
  \item El programa de fidelización con tarjeta digital y canjes
  \item El backoffice de administración (usuarios, cafés, menú, reseñas, animales)
  \item Los paneles de operación: recepción, cocina (\gls{kds}), manager y keeper
  \item La \gls{api} \gls{rest} para integración con clientes móviles y terceros
\end{enumerate}

\section{Stack tecnológico}
\label{sec:stack}

La \cref{tab:stack} resume las tecnologías principales utilizadas en el proyecto.

\begin{longtable}{@{}lll@{}}
  \caption{Stack tecnológico del proyecto}
  \label{tab:stack}\\
  \toprule
  \textbf{Capa} & \textbf{Tecnología} & \textbf{Versión / Detalles}\\
  \midrule
  \endfirsthead
  \toprule
  \textbf{Capa} & \textbf{Tecnología} & \textbf{Versión / Detalles}\\
  \midrule
  \endhead
  \bottomrule
  \endfoot
  Lenguaje backend   & PHP                  & 8.4 (strict types en todos los archivos)\\
  Framework          & MVC custom           & Sin Laravel / Symfony, PSR-7/14/15\\
  Servidor HTTP      & FrankenPHP           & Caddy + PHP integrado\\
  Base de datos      & MySQL                & 8.0, 19 migraciones SQL\\
  Caché / Cola       & Redis                & 7.x, sesiones + cola de trabajos\\
  Frontend           & Alpine.js            & 3.x, sin bundler en producción\\
  CSS                & CSS3 custom          & Variables, Grid, Flexbox\\
  Contenedores       & Docker Compose       & Multistage build\\
  Despliegue         & Railway              & + GitHub Actions \gls{ci}\\
  Tests unitarios    & PHPUnit              & 13.x, tests paralelos\\
  Tests E2E          & Playwright           & TypeScript\\
  Análisis estático  & PHPStan + Psalm      & Nivel 5 + baseline\\
  Email              & Mailpit (dev)        & SMTP local\\
  Facturas PDF       & Dompdf               & Generación asíncrona\\
\end{longtable}

\section{Metodología de desarrollo}
\label{sec:metodologia}

El desarrollo siguió una metodología ágil inspirada en \textbf{Scrum}, adaptada al
contexto de un equipo unipersonal. Se utilizó Trello como tablero Kanban con las
columnas \textit{Backlog}, \textit{En progreso}, \textit{Revisión} y \textit{Hecho}.

El control de versiones se gestionó con Git siguiendo el flujo \textbf{Git Flow}:
\begin{itemize}[nosep]
  \item \texttt{main} — código en producción (Railway)
  \item \texttt{develop} — integración continua
  \item \texttt{feature/\textit{nombre}} — desarrollo de funcionalidades
  \item \texttt{hotfix/\textit{nombre}} — correcciones urgentes
\end{itemize}

Cada \textit{pull request} requería superar el \textit{quality gate} completo:
PHPStan nivel~5, Psalm, suite PHPUnit y verificación de estilo PSR-12.

% ============================================================
%  CAPÍTULO 2 — DESARROLLO TÉCNICO / MANUAL TÉCNICO
% ============================================================
\chapter{Desarrollo Técnico}
\label{chap:tecnico}

\section{Arquitectura del sistema}
\label{sec:arquitectura}

Komorebi Café implementa el patrón \gls{mvc} desde cero, sin delegar en ningún
framework de aplicación externo. El objetivo es demostrar comprensión profunda de los
mecanismos subyacentes: enrutamiento, inyección de dependencias, ciclo de
petición-respuesta y gestión de middleware.

\subsection{Capas principales}

\begin{description}[style=nextline]
  \item[\texttt{public/index.php} — Front controller]
    Único punto de entrada. Lee variables de entorno, inicializa el contenedor de
    dependencias y pasa la petición al router.
  \item[\texttt{bootstrap/container.php} — Contenedor IoC]
    Registra todos los \textit{Service Providers}. Ciclo \texttt{register→boot}.
    Singletons enlazados por interfaz.
  \item[\texttt{app/routes.php} — Definición de rutas]
    250~rutas agrupadas por dominio con prefijos y middleware declarados de
    forma explícita.
  \item[\texttt{app/Core/} — Framework interno]
    Router PSR-7/15, Contenedor IoC, Vista (auto-escape + Raw), Base de datos (PDO),
    Caché (Redis + local), Cola (Redis), Logger (stdout 12-Factor).
  \item[\texttt{app/Http/Controllers/} — Controladores]
    Agrupados en: \texttt{Admin/}, \texttt{Api/V1/}, \texttt{Auth/}, \texttt{Kitchen/},
    \texttt{Keeper/}, \texttt{Manager/}, \texttt{Public/}, \texttt{Reception/},
    \texttt{Shared/}, \texttt{Supervisor/}, \texttt{User/}.
  \item[\texttt{app/Services/} — Lógica de negocio]
    53 servicios. Todos devuelven \texttt{Result::ok()/fail()} — nunca lanzan
    excepciones para fallos esperados.
  \item[\texttt{app/Repositories/} — Capa de datos]
    29 repositorios. Extienden \texttt{AbstractRepository} con PDO directo.
  \item[\texttt{app/Events/ + Listeners/} — Eventos asíncronos]
    PSR-14 con Symfony EventDispatcher. Procesados por workers Redis.
  \item[\texttt{migrations/} — Esquema de base de datos]
    19 archivos SQL numerados, aplicables mediante \texttt{make db-migrate}.
\end{description}

\subsection{Diagrama de contexto C4 (nivel 1)}

\begin{figure}[H]
  \centering
  \begin{tikzpicture}[
    box/.style={draw,rounded corners,fill=komorebi-light,
                text width=3cm,align=center,minimum height=1.2cm,
                font=\small},
    ext/.style={draw,rounded corners,fill=gray!15,
                text width=3cm,align=center,minimum height=1.2cm,
                font=\small},
    arr/.style={-{Latex},thick}
  ]
    \node[box] (app) {\textbf{Komorebi Café}\\Web App};
    \node[ext,left=2.5cm of app] (cliente) {Cliente\\(navegador)};
    \node[ext,right=2.5cm of app] (railway) {Railway\\(producción)};
    \node[ext,above=1.8cm of app] (email) {Mailpit /\\SMTP externo};
    \node[ext,below=1.8cm of app] (mysql) {MySQL 8.0\\+ Redis 7};

    \draw[arr] (cliente) -- node[above,font=\scriptsize]{HTTPS} (app);
    \draw[arr] (app) -- node[above,font=\scriptsize]{Deploy} (railway);
    \draw[arr] (app) -- node[right,font=\scriptsize]{SMTP} (email);
    \draw[arr] (app) -- node[right,font=\scriptsize]{PDO / Redis} (mysql);
  \end{tikzpicture}
  \caption{Diagrama de contexto C4 — nivel sistema}
  \label{fig:c4-contexto}
\end{figure}

\section{Ciclo de petición HTTP y pipeline de middleware}
\label{sec:ciclo-http}

Cada petición HTTP atraviesa un pipeline de quince middlewares \gls{psr}-15 antes de
llegar al controlador correspondiente. La \cref{fig:pipeline} muestra el orden de
ejecución.

\begin{figure}[H]
  \centering
  \begin{tikzpicture}[
    mw/.style={draw,fill=komorebi-light,rounded corners,
               text width=3.2cm,align=center,minimum height=0.7cm,
               font=\scriptsize},
    arr/.style={-{Latex},thick,komorebi}
  ]
    \foreach \label/\y in {
      ErrorHandler/0,
      SecurityHeaders/-1,
      RequestLog/-2,
      Session/-3,
      CSRF/-4,
      RateLimit/-5,
      PayloadSize/-6,
      CORS/-7,
      Auth/-8,
      ApiAuth/-9,
      Role/-10,
      ApiRole/-11,
      Authorization/-12,
      CafeScope/-13,
      Idempotency/-14
    }{
      \node[mw] (\label) at (0,\y) {\label};
    }
    \foreach \a/\b in {
      ErrorHandler/SecurityHeaders,
      SecurityHeaders/RequestLog,
      RequestLog/Session,
      Session/CSRF,
      CSRF/RateLimit,
      RateLimit/PayloadSize,
      PayloadSize/CORS,
      CORS/Auth,
      Auth/ApiAuth,
      ApiAuth/Role,
      Role/ApiRole,
      ApiRole/Authorization,
      Authorization/CafeScope,
      CafeScope/Idempotency
    }{
      \draw[arr] (\a) -- (\b);
    }
    \node[right=0.5cm of ErrorHandler,font=\scriptsize,text width=5cm]
      {Captura toda excepción; devuelve página de error apropiada};
    \node[right=0.5cm of Session,font=\scriptsize,text width=5cm]
      {Inicia sesión PHP + regeneración periódica};
    \node[right=0.5cm of CSRF,font=\scriptsize,text width=5cm]
      {Valida token en POST/PUT/PATCH/DELETE web};
    \node[right=0.5cm of Auth,font=\scriptsize,text width=5cm]
      {Verifica sesión autenticada (rutas web protegidas)};
    \node[right=0.5cm of ApiAuth,font=\scriptsize,text width=5cm]
      {Bearer token → hash SHA-256 → \texttt{api\_tokens}};
    \node[right=0.5cm of Idempotency,font=\scriptsize,text width=5cm]
      {UUID v4 header → Redis 24h (solo POST /api/v1/reservations)};
  \end{tikzpicture}
  \caption{Pipeline de 15 middlewares PSR-15}
  \label{fig:pipeline}
\end{figure}

La \cref{tab:middlewares} describe brevemente cada middleware.

\begin{longtable}{@{}lp{7cm}@{}}
  \caption{Catálogo de middlewares PSR-15}
  \label{tab:middlewares}\\
  \toprule
  \textbf{Middleware} & \textbf{Responsabilidad}\\
  \midrule
  \endfirsthead
  \toprule
  \textbf{Middleware} & \textbf{Responsabilidad}\\
  \midrule
  \endhead
  \bottomrule
  \endfoot
  \texttt{ErrorHandlerMiddleware}
    & Envuelve toda la cadena; captura \texttt{Throwable} y devuelve respuesta
      de error apropiada (HTML 4xx/5xx o JSON según el tipo de ruta).\\
  \texttt{SecurityHeadersMiddleware}
    & Añade cabeceras de seguridad: \texttt{X-Frame-Options}, \texttt{X-Content-Type-Options},
      \texttt{Content-Security-Policy}, \texttt{Referrer-Policy}.\\
  \texttt{RequestLogMiddleware}
    & Registra método, URI, IP, agente y tiempo de respuesta en stdout (12-Factor XI).\\
  \texttt{SessionMiddleware}
    & Inicia la sesión PHP, gestiona regeneración periódica del identificador y
      establece atributos de sesión en el objeto \texttt{Request}.\\
  \texttt{CsrfMiddleware}
    & Verifica el token CSRF en todas las rutas mutantes (POST/PUT/PATCH/DELETE) web.
      Rechaza con 419 si el token falta o no coincide.\\
  \texttt{RateLimitMiddleware}
    & Limitación de velocidad por clave (bucket Redis): \texttt{login},
      \texttt{registration}, \texttt{password\_reset}. Responde 429 al exceder.\\
  \texttt{PayloadSizeMiddleware}
    & Rechaza cuerpos de petición superiores a 10~MB para prevenir ataques de
      agotamiento de memoria.\\
  \texttt{CorsMiddleware}
    & Gestiona cabeceras CORS y respuestas \texttt{OPTIONS} preflight para la API.\\
  \texttt{AuthMiddleware}
    & Requiere sesión autenticada. Redirige a \texttt{/login} si no está logado.\\
  \texttt{ApiAuthMiddleware}
    & Extrae el Bearer token, aplica SHA-256, consulta \texttt{api\_tokens} y
      autentica la petición de forma stateless.\\
  \texttt{RoleMiddleware}
    & Verifica que el rol del usuario en sesión coincide con el rol requerido
      por la ruta web.\\
  \texttt{ApiRoleMiddleware}
    & Igual que \texttt{RoleMiddleware} pero para peticiones API autenticadas
      con Bearer token.\\
  \texttt{AuthorizationMiddleware}
    & Aplica políticas de autorización de grano fino (p.ej.\ solo el propietario
      puede modificar sus datos).\\
  \texttt{CafeScopeMiddleware}
    & Restringe el acceso de managers y supervisors a los datos del café al que
      están asignados.\\
  \texttt{IdempotencyMiddleware}
    & Verifica el header \texttt{Idempotency-Key} (UUID~v4); devuelve la respuesta
      cacheada (Redis, TTL 24~h) si existe, o ejecuta y almacena el resultado.\\
\end{longtable}

\section{Flujos de autenticación}
\label{sec:auth}

El sistema implementa tres mecanismos de autenticación diferenciados según el tipo
de cliente.

\subsection{Login web con sesión}

\begin{figure}[H]
  \centering
  \begin{sequencediagram}
    \newthread{U}{Usuario}
    \newthread{C}{AuthController}
    \newthread{S}{AuthService}
    \newthread{R}{UserRepository}
    \newthread{Se}{Session}

    \begin{call}{U}{POST /login (email, password, csrf\_token)}{C}{}
      \begin{call}{C}{processLogin(\$request)}{S}{}
        \begin{call}{S}{findByEmail(\$email)}{R}{User|null}
        \end{call}
        \begin{call}{S}{password\_verify(\$plain, \$hash)}{S}{bool}
        \end{call}
      \end{call}
      \begin{call}{C}{set(\$userId, \$role)}{Se}{}
      \end{call}
    \end{call}
  \end{sequencediagram}
  \caption{Diagrama de secuencia: login web}
  \label{fig:seq-login}
\end{figure}

\subsection{Verificación \gls{rbac} por middleware}

El \texttt{RoleMiddleware} comprueba el rol almacenado en sesión antes de que la
petición llegue al controlador. Los roles disponibles son:
\texttt{admin}, \texttt{manager}, \texttt{supervisor}, \texttt{reception},
\texttt{kitchen}, \texttt{keeper} y \texttt{user}.

\subsection{API stateless con Bearer token}

\begin{figure}[H]
  \centering
  \begin{sequencediagram}
    \newthread{C}{Cliente API}
    \newthread{TC}{TokenController}
    \newthread{TS}{ApiTokenService}
    \newthread{R}{UserRepository}
    \newthread{DB}{api\_tokens}

    \begin{call}{C}{POST /api/v1/tokens (email, password)}{TC}{}
      \begin{call}{TC}{create(\$userId, \$name, \$abilities)}{TS}{}
        \begin{call}{TS}{\texttt{random\_bytes(32)} → plain token}{TS}{}
        \end{call}
        \begin{call}{TS}{\texttt{hash('sha256', token)}}{TS}{hash}
        \end{call}
        \begin{call}{TS}{insert(token\_hash, expires\_at, abilities)}{DB}{}
        \end{call}
      \end{call}
    \end{call}
    \mess{TC}{plain token (devuelto una sola vez)}{C}

    \begin{sdblock}{Petición posterior}{}
      \begin{call}{C}{GET /api/v1/... Bearer <token>}{TC}{}
        \begin{call}{TC}{\texttt{hash('sha256', token)}}{TC}{hash}
        \end{call}
        \begin{call}{TC}{findByHash(hash)}{DB}{ApiToken}
        \end{call}
      \end{call}
    \end{sdblock}
  \end{sequencediagram}
  \caption{Diagrama de secuencia: autenticación Bearer token API}
  \label{fig:seq-token}
\end{figure}

\section{Flujo de reservas con control de idempotencia}
\label{sec:reservas-flujo}

\begin{figure}[H]
  \centering
  \begin{sequencediagram}
    \newthread{U}{Usuario}
    \newthread{IM}{IdempotencyMW}
    \newthread{RC}{ReservationController}
    \newthread{RS}{ReservationService}
    \newthread{DB}{reservations}
    \newthread{RD}{Redis}
    \newthread{EV}{EventDispatcher}

    \begin{call}{U}{POST /api/v1/reservations\\Idempotency-Key: UUID v4}{IM}{}
      \begin{call}{IM}{GET idempotency:\{key\}}{RD}{hit/miss}
      \end{call}
      \begin{sdblock}{Cache HIT}{}
        \mess{IM}{Response cacheada (200/201)}{U}
      \end{sdblock}
      \begin{sdblock}{Cache MISS}{}
        \begin{call}{IM}{handle(\$request)}{RC}{}
          \begin{call}{RC}{create(\$dto)}{RS}{}
            \begin{call}{RS}{INSERT}{DB}{Reservation}
            \end{call}
            \begin{call}{RS}{dispatch(ReservationCreatedEvent)}{EV}{}
            \end{call}
          \end{call}
        \end{call}
        \begin{call}{IM}{SET idempotency:\{key\} TTL 86400}{RD}{}
        \end{call}
        \mess{IM}{201 Created (invoice\_pdf\_url: null)}{U}
      \end{sdblock}
    \end{call}
  \end{sequencediagram}
  \caption{Diagrama de secuencia: creación de reserva con idempotencia}
  \label{fig:seq-reservas}
\end{figure}

\begin{infobox}
  El campo \texttt{invoice\_pdf\_url} es \texttt{NULL} en la respuesta inicial.
  El \texttt{InvoicePDFService} genera el PDF de forma asíncrona (worker Redis) y
  actualiza la columna \texttt{reservations.invoice\_pdf\_url} cuando el archivo
  está disponible.
\end{infobox}

\section{Modelo entidad-relación}
\label{sec:er}

La base de datos está compuesta por 19 migraciones SQL que definen las tablas
principales del sistema. La \cref{tab:entidades} resume las entidades más relevantes.

\begin{longtable}{@{}lp{5cm}p{5cm}@{}}
  \caption{Entidades principales del modelo de datos}
  \label{tab:entidades}\\
  \toprule
  \textbf{Tabla} & \textbf{Propósito} & \textbf{Columnas clave}\\
  \midrule
  \endfirsthead
  \toprule
  \textbf{Tabla} & \textbf{Propósito} & \textbf{Columnas clave}\\
  \midrule
  \endhead
  \bottomrule
  \endfoot
  \texttt{users}
    & Cuentas de usuario
    & \texttt{id, email, password\_hash, role, is\_active}\\
  \texttt{cafes}
    & Establecimientos temáticos
    & \texttt{id, slug, name, capacity, is\_active}\\
  \texttt{products}
    & Carta del menú
    & \texttt{id, cafe\_id, name, price, is\_available}\\
  \texttt{allergens}
    & Alérgenos por producto
    & \texttt{id, product\_id, name, severity}\\
  \texttt{reservations}
    & Reservas de clientes
    & \texttt{id, cafe\_id, user\_id, date, guests,}\\
    & & \texttt{status, idempotency\_key, invoice\_pdf\_url}\\
  \texttt{api\_tokens}
    & Tokens de acceso API
    & \texttt{id, user\_id, token\_hash, abilities,}\\
    & & \texttt{expires\_at, revoked\_at}\\
  \texttt{loyalty\_cards}
    & Tarjetas de fidelización
    & \texttt{id, user\_id, points, tier}\\
  \texttt{animals}
    & Animales de los cafés
    & \texttt{id, cafe\_id, name, species, is\_active}\\
  \texttt{animal\_incidents}
    & Incidencias de animales
    & \texttt{id, animal\_id, severity, description, resolution}\\
  \texttt{health\_checks}
    & Revisiones de salud
    & \texttt{id, animal\_id, keeper\_id, checked\_at, notes}\\
  \texttt{reviews}
    & Reseñas de clientes
    & \texttt{id, cafe\_id, user\_id, rating, body, status}\\
  \texttt{newsletter\_subscriptions}
    & Suscriptores de newsletter
    & \texttt{id, email, verified\_at, expires\_at}\\
  \texttt{staff\_shifts}
    & Turnos de personal
    & \texttt{id, user\_id, cafe\_id, start\_time, end\_time}\\
  \texttt{time\_slots}
    & Franjas horarias de reserva
    & \texttt{id, cafe\_id, date, time, capacity, booked}\\
  \texttt{waitlist}
    & Lista de espera
    & \texttt{id, cafe\_id, user\_id, token, status}\\
\end{longtable}

El diagrama entidad-relación completo (en formato PlantUML) se encuentra en
\texttt{docs/diagrams/er.puml} e incluye todas las relaciones de clave foránea
y constraints de unicidad.

\section{Implementación frontend}
\label{sec:frontend}

Las vistas se renderizan en el servidor mediante el motor de plantillas PHP interno
(\texttt{app/Core/View.php}). Alpine.js proporciona la interactividad del lado del
cliente sin necesidad de un \textit{bundler} en producción.

\subsection{Convenciones de vistas}
\begin{itemize}[nosep]
  \item Todas las vistas viven en \texttt{resources/views/}.
  \item \texttt{View::render()} auto-escapa todos los valores de \texttt{\$data}.
  \item Para HTML o JSON pre-procesados se usa \texttt{Raw::html()} / \texttt{Raw::json()}.
  \item Los layouts se encuentran en \texttt{resources/views/layouts/}:
    \texttt{main.php} (público/usuario), \texttt{backoffice.php} (admin/manager),
    \texttt{kds.php} (cocina), \texttt{mobile.php}, \texttt{reception.php}.
\end{itemize}

\subsection{Patrón de datos para Alpine.js}

\begin{phpcode}{Pasar JSON seguro a Alpine.js}
<?php declare(strict_types=1);

// En el controlador:
$data = [
    'reservations' => Raw::json($reservationsArray),
    'cafes'        => Raw::json($cafesArray),
];
View::render('admin/reservations/index', $data, ['admin-reservations.css']);
\end{phpcode}

\section{Implementación backend}
\label{sec:backend}

\subsection{El Result pattern}

Todos los servicios devuelven un objeto \texttt{Result} en lugar de lanzar excepciones
para fallos esperados. Esto hace que el flujo de control sea explícito y predecible.

\begin{phpcode}{Ejemplo de uso del Result pattern}
<?php declare(strict_types=1);

namespace App\Services;

use App\Core\Result;

final class ReservationService
{
    public function create(ReservationDTO $dto): Result
    {
        if ($dto->guests > $this->getCapacity($dto->cafeId)) {
            return Result::fail('Capacidad insuficiente', 'capacity_exceeded');
        }

        $id = $this->repository->insert($dto->toArray());
        return Result::ok(['id' => $id]);
    }
}
\end{phpcode}

\begin{phpcode}{Consumo del Result en el controlador}
<?php declare(strict_types=1);

$result = $this->service->create($dto);

if (!$result->ok) {
    Flash::error($result->getMessage());
    return $this->response->redirect('/reservar');
}

return $this->response->redirect('/reservas/confirmacion/' . $result->data['id']);
\end{phpcode}

\subsection{Catálogo de rutas}
\label{subsec:rutas}

\begin{landscape}
\begin{longtable}{@{}llll@{}}
  \caption{Catálogo de rutas del sistema (selección)}
  \label{tab:rutas}\\
  \toprule
  \textbf{Método} & \textbf{Ruta} & \textbf{Controlador} & \textbf{Middleware}\\
  \midrule
  \endfirsthead
  \toprule
  \textbf{Método} & \textbf{Ruta} & \textbf{Controlador} & \textbf{Middleware}\\
  \midrule
  \endhead
  \bottomrule
  \endfoot
  % ── Públicas ──
  GET  & \texttt{/}                    & Public\textbackslash HomeController@index          & —\\
  GET  & \texttt{/cafes}               & Public\textbackslash CafeController@index          & —\\
  GET  & \texttt{/cafes/\{slug\}}      & Public\textbackslash CafeController@show           & —\\
  GET  & \texttt{/menu}                & Public\textbackslash MenuController@index          & —\\
  GET  & \texttt{/quiz}                & Public\textbackslash QuizController@index          & —\\
  POST & \texttt{/quiz/resultado}      & Public\textbackslash QuizController@resultado      & csrf\\
  GET  & \texttt{/historia}            & Public\textbackslash PageController@historia       & —\\
  GET  & \texttt{/faq}                 & Public\textbackslash PageController@faq            & —\\
  GET  & \texttt{/contacto}            & Public\textbackslash PageController@contacto       & —\\
  % ── Autenticación ──
  GET  & \texttt{/login}               & Auth\textbackslash AuthController@showLogin         & guest\\
  POST & \texttt{/login}               & Auth\textbackslash AuthController@processLogin      & csrf, rateLimit(login)\\
  GET  & \texttt{/registro}            & Auth\textbackslash AuthController@showRegister      & guest\\
  POST & \texttt{/registro}            & Auth\textbackslash AuthController@processRegister   & csrf, rateLimit(registration)\\
  GET  & \texttt{/forgot-password}     & Auth\textbackslash PasswordResetController         & guest\\
  POST & \texttt{/forgot-password}     & Auth\textbackslash PasswordResetController         & csrf, rateLimit(password\_reset)\\
  POST & \texttt{/logout}              & Auth\textbackslash AuthController@logout            & csrf\\
  % ── Usuario autenticado ──
  GET  & \texttt{/perfil}              & Shared\textbackslash UserController@profile        & auth\\
  GET  & \texttt{/reservas/mis-reservas}& Shared\textbackslash ReservationController        & auth\\
  GET  & \texttt{/loyalty/card}        & Public\textbackslash LoyaltyController@card        & auth\\
  GET  & \texttt{/mis-favoritos}       & User\textbackslash FavoriteController@index        & auth\\
  % ── Admin ──
  GET  & \texttt{/admin/dashboard}     & Admin\textbackslash DashboardController@index      & auth, role(admin)\\
  GET  & \texttt{/admin/users}         & Admin\textbackslash UserController@index           & auth, role(admin)\\
  GET  & \texttt{/admin/cafes}         & Admin\textbackslash CafeController@index           & auth, role(admin)\\
  GET  & \texttt{/admin/menu}          & Admin\textbackslash MenuController@index           & auth, role(admin)\\
  GET  & \texttt{/admin/reservations}  & Admin\textbackslash ReservationController@index    & auth, role(admin)\\
  GET  & \texttt{/admin/animals}       & Admin\textbackslash AnimalController@index         & auth, role(admin)\\
  GET  & \texttt{/admin/reports}       & Admin\textbackslash ReportController@index         & auth, role(admin)\\
  GET  & \texttt{/admin/settings}      & Admin\textbackslash SystemController@settings      & auth, role(admin)\\
  GET  & \texttt{/admin/logs/audit}    & Admin\textbackslash AuditLogController@index       & auth, role(admin)\\
  GET  & \texttt{/admin/loyalty}       & Admin\textbackslash LoyaltyController@index        & auth, role(admin)\\
  % ── Manager ──
  GET  & \texttt{/manager/dashboard}   & Manager\textbackslash DashboardController@index   & auth, role(manager)\\
  GET  & \texttt{/manager/products}    & Manager\textbackslash ProductController@index      & auth, role(manager), ownsCafe\\
  GET  & \texttt{/manager/staff}       & Manager\textbackslash StaffController@index        & auth, role(manager), ownsCafe\\
  % ── Recepción / Cocina ──
  GET  & \texttt{/reception}           & Reception\textbackslash ReceptionController@index  & auth, role(reception)\\
  GET  & \texttt{/kitchen}             & Kitchen\textbackslash KitchenController@index      & auth, role(kitchen)\\
  % ── Keeper ──
  GET  & \texttt{/keeper/dashboard}    & Keeper\textbackslash AnimalDashboardController     & auth, role(keeper)\\
  GET  & \texttt{/keeper/health-checks}& Keeper\textbackslash HealthCheckController@index  & auth, role(keeper)\\
  % ── API v1 pública ──
  GET  & \texttt{/api/v1/menu/alergenos}& Api\textbackslash V1\textbackslash MenuController@allergens & cors\\
  GET  & \texttt{/api/v1/cafes}        & Api\textbackslash V1\textbackslash CafeController@index    & cors\\
  GET  & \texttt{/api/v1/time-slots/available}& Api\textbackslash V1\textbackslash TimeSlotController & cors\\
  % ── API v1 autenticada ──
  POST & \texttt{/api/v1/reservations} & Api\textbackslash V1\textbackslash ReservationController@create & cors, apiAuth, csrf, idempotency\\
  GET  & \texttt{/api/v1/tokens}       & Api\textbackslash V1\textbackslash TokenController@list        & cors, apiAuth\\
  POST & \texttt{/api/v1/tokens}       & Api\textbackslash V1\textbackslash TokenController@create      & cors, csrf\\
  DELETE& \texttt{/api/v1/tokens/\{id\}}& Api\textbackslash V1\textbackslash TokenController@revoke     & cors, apiAuth, csrf\\
  % ── API Admin ──
  POST & \texttt{/api/v1/admin/users}  & Api\textbackslash V1\textbackslash Admin\textbackslash UserApiController@create & apiAuth, role(admin), csrf\\
  PUT  & \texttt{/api/v1/admin/cafes/\{id\}}& Api\textbackslash V1\textbackslash Admin\textbackslash CafeApiController@update & apiAuth, role(admin), csrf\\
  POST & \texttt{/api/v1/admin/menu}   & Api\textbackslash V1\textbackslash Admin\textbackslash MenuApiController@create & apiAuth, role(admin), csrf\\
  % ── API Operaciones ──
  POST & \texttt{/api/v1/ops/reception/reservations/\{id\}/checkin}& Api\textbackslash V1\textbackslash Ops\textbackslash ReceptionApiController & apiAuth, role(reception), csrf\\
  POST & \texttt{/api/v1/ops/kitchen/orders/\{id\}/complete}& Api\textbackslash V1\textbackslash Ops\textbackslash KitchenApiController & apiAuth, role(kitchen), csrf\\
\end{longtable}
\end{landscape}

\subsection{Catálogo de servicios de negocio}
\label{subsec:servicios}

\begin{longtable}{@{}lp{8cm}@{}}
  \caption{Catálogo de servicios de negocio (53 servicios)}
  \label{tab:servicios}\\
  \toprule
  \textbf{Servicio} & \textbf{Responsabilidad principal}\\
  \midrule
  \endfirsthead
  \toprule
  \textbf{Servicio} & \textbf{Responsabilidad principal}\\
  \midrule
  \endhead
  \bottomrule
  \endfoot
  \texttt{AuthService}                 & Autenticación web: login, logout, verificación de credenciales\\
  \texttt{ApiTokenService}             & Gestión de tokens Bearer: creación, revocación, validación\\
  \texttt{AuthTokenService}            & Tokens de verificación de email y reset de contraseña\\
  \texttt{PasswordResetService}        & Flujo completo de recuperación de contraseña\\
  \texttt{EmailVerificationService}    & Envío y verificación de emails de confirmación de cuenta\\
  \texttt{SessionManagementService}    & Listado, revocación y expiración de sesiones activas\\
  \texttt{UserAccountService}          & Alta, baja y gestión general de cuentas de usuario\\
  \texttt{UserManagementService}       & Operaciones admin sobre usuarios (toggle activo, roles)\\
  \texttt{UserProfileService}          & Actualización de perfil, avatar y preferencias de usuario\\
  \texttt{UserPreferenceService}       & Preferencias de cookies, filtros dietéticos y recientes\\
  \texttt{CafeService}                 & CRUD de cafés temáticos y gestión de estado\\
  \texttt{MenuService}                 & Catálogo de menú: categorías y disponibilidad\\
  \texttt{ProductService}              & CRUD de productos de carta con toggle de disponibilidad\\
  \texttt{AllergenService}             & Asociación de alérgenos a productos\\
  \texttt{ReservationService}          & Creación, confirmación y cancelación de reservas\\
  \texttt{ReservationTimeSlotService}  & Validación y gestión de franjas horarias de reserva\\
  \texttt{AvailabilityService}         & Cálculo de plazas disponibles por café, fecha y hora\\
  \texttt{TimeSlotService}             & CRUD de time slots y estadísticas de ocupación\\
  \texttt{WaitlistService}             & Alta en lista de espera, posición y confirmación\\
  \texttt{LoyaltyService}              & Acumulación de puntos, canjes y gestión de catálogo\\
  \texttt{GamificationService}         & Logros, niveles y recompensas del programa de fidelización\\
  \texttt{ReviewService}               & Creación y gestión de reseñas de clientes\\
  \texttt{ReviewModerationService}     & Aprobación y rechazo de reseñas por moderadores\\
  \texttt{ReviewQueryService}          & Consultas y filtros avanzados sobre reseñas\\
  \texttt{CartService}                 & Gestión del carrito (Redis, por sesión o usuario)\\
  \texttt{RecentlyViewedService}       & Historial de productos visitados recientemente (cookies)\\
  \texttt{NewsletterService}           & Suscripción, verificación y baja del boletín\\
  \texttt{EmailService}                & Envío de emails transaccionales (queue + SMTP)\\
  \texttt{InvoicePDFService}           & Generación asíncrona de facturas PDF (Dompdf)\\
  \texttt{KitchenService}              & Órdenes activas e historial del KDS de cocina\\
  \texttt{ReceptionService}            & Check-in, checkout y procesamiento de pagos en recepción\\
  \texttt{AnimalCareService}           & CRUD de animales, incidentes y revisiones de salud\\
  \texttt{HealthCheckService}          & Registro y consulta de revisiones veterinarias\\
  \texttt{StaffShiftService}           & Gestión de turnos de personal por café\\
  \texttt{SupervisorAssignmentService} & Asignaciones de supervisor a café y turnos\\
  \texttt{SettingsService}             & Lectura y escritura de configuración del sistema\\
  \texttt{CacheService}                & Abstracción de Redis con tags e invalidación por grupo\\
  \texttt{RateLimitingService}         & Implementación de ventana deslizante en Redis\\
  \texttt{FileUploadService}           & Subida y validación de imágenes (avatares, animales)\\
  \texttt{CloudinaryStorageService}    & Almacenamiento de imágenes en Cloudinary (producción)\\
  \texttt{TelegramService}             & Notificaciones push al bot de Telegram del equipo\\
  \texttt{MercurePublisherService}     & Publicación de eventos SSE en el hub Mercure (KDS)\\
  \texttt{WeatherService}              & Consulta de meteorología local (API externa)\\
  \texttt{NavigationService}           & Generación del menú de navegación según rol activo\\
  \texttt{AdminActivityService}        & Registro de acciones de administración en audit log\\
  \texttt{AdminStatisticsService}      & Cálculo de KPIs para el dashboard de administración\\
  \texttt{AdminReportService}          & Generación y exportación de informes en CSV/PDF\\
  \texttt{AccountDeletionService}      & Eliminación segura de cuentas de usuario y datos asociados\\
  \texttt{ClimaContextoService}        & Contexto climático (temperatura, lluvia) para personalizar la experiencia\\
  \texttt{ContextServiceInstance}      & Versión inyectable de ContextService: resuelve el café activo por sesión\\
  \texttt{DashboardService}            & Métricas en tiempo real para el panel de control del gestor\\
  \texttt{FestivosJaponesesService}    & Consulta de festivos japoneses (祝日 Shukujitsu) para validación de slots\\
  \texttt{MicroestacionesService}      & Calendario solar japonés (24 Sekki) para personalización estacional\\
\end{longtable}

\subsection{Catálogo de repositorios}
\label{subsec:repositorios}

\begin{longtable}{@{}lp{7cm}@{}}
  \caption{Catálogo de repositorios (29 repositorios)}
  \label{tab:repos}\\
  \toprule
  \textbf{Repositorio} & \textbf{Entidad / Tabla gestionada}\\
  \midrule
  \endfirsthead
  \toprule
  \textbf{Repositorio} & \textbf{Entidad / Tabla gestionada}\\
  \midrule
  \endhead
  \bottomrule
  \endfoot
  \texttt{UserRepository}                  & \texttt{users}\\
  \texttt{RoleRepository}                  & \texttt{roles}, \texttt{user\_roles}\\
  \texttt{SessionRepository}               & \texttt{sessions}\\
  \texttt{AuthTokenRepository}             & \texttt{auth\_tokens} (email/reset)\\
  \texttt{ApiTokenRepository}              & \texttt{api\_tokens} (Bearer)\\
  \texttt{CafeRepository}                  & \texttt{cafes}\\
  \texttt{MenuCategoryRepository}          & \texttt{menu\_categories}\\
  \texttt{MenuRepository}                  & \texttt{menu\_items}\\
  \texttt{ProductRepository}               & \texttt{products}\\
  \texttt{AllergenRepository}              & \texttt{allergens}\\
  \texttt{ReservationRepository}           & \texttt{reservations}\\
  \texttt{ReservationItemRepository}       & \texttt{reservation\_items}\\
  \texttt{TimeSlotRepository}              & \texttt{time\_slots}\\
  \texttt{WaitlistRepository}              & \texttt{waitlist}\\
  \texttt{ReviewRepository}                & \texttt{reviews}\\
  \texttt{LoyaltyRepository}               & \texttt{loyalty\_cards}, \texttt{loyalty\_transactions}\\
  \texttt{FavoriteRepository}              & \texttt{favorites}\\
  \texttt{NewsletterSubscriptionRepository}& \texttt{newsletter\_subscriptions}\\
  \texttt{AnimalRepository}                & \texttt{animals}\\
  \texttt{AnimalIncidentRepository}        & \texttt{animal\_incidents}\\
  \texttt{HealthCheckRepository}           & \texttt{health\_checks}\\
  \texttt{StaffShiftRepository}            & \texttt{staff\_shifts}\\
  \texttt{SupervisorAssignmentRepository}  & \texttt{supervisor\_assignments}\\
  \texttt{SettingRepository}               & \texttt{system\_settings}\\
  \texttt{AuditLogRepository}              & \texttt{audit\_logs}\\
  \texttt{AuthLogRepository}               & \texttt{auth\_logs}\\
  \texttt{StatisticsRepository}            & Vistas agregadas / consultas analytics\\
  \texttt{TrackerRepository}               & \texttt{interaction\_sessions}\\
  \texttt{InteractionSessionRepository}    & Sesiones de interacción (A/B testing)\\
\end{longtable}

\section{Pruebas}
\label{sec:pruebas}

\subsection{Tests unitarios con PHPUnit}

Los tests unitarios se ubican en \texttt{tests/Unit/} replicando la estructura de
\texttt{app/}. Cada archivo de test incluye el docblock obligatorio del proyecto:

\begin{phpcode}{Estructura de test unitario}
<?php declare(strict_types=1);

/**
 * ¿Qué pruebas aquí?
 *    El ReservationService crea reservas correctamente y falla
 *    con gracia cuando se supera la capacidad.
 * ¿Qué me quieres demostrar?
 *    Que el Result pattern encapsula los errores sin excepciones.
 * ¿Qué va a fallar en este test si se cambia el código?
 *    Si create() lanza una excepción en lugar de devolver Result::fail().
 */
class ReservationServiceTest extends TestCase
{
    public function testCreateReturnsOkWhenCapacityIsAvailable(): void
    {
        $repo = $this->createStub(ReservationRepository::class);
        $repo->method('insert')->willReturn(42);

        $service = new ReservationService($repo, /* ... */);
        $result  = $service->create($dto);

        $this->assertTrue($result->ok);
        $this->assertSame(42, $result->data['id']);
    }
}
\end{phpcode}

\subsection{Tests de integración}

Los tests de integración en \texttt{tests/Integration/} utilizan una base de datos
real (\texttt{docker-compose.test.yml}) con migraciones y seeders ejecutados antes
de cada suite. Se ejecutan con:

\begin{bashcode}{Ejecutar tests de integración}
make test-integration
\end{bashcode}

\subsection{Tests E2E con Playwright}

Los tests end-to-end se encuentran en \texttt{tests/e2e/} y están escritos en
TypeScript. Cubren flujos completos: login, creación de reserva, check-in y
verificación de KDS. Se ejecutan con:

\begin{bashcode}{Ejecutar tests E2E}
make e2e
make e2e-a11y   # solo accesibilidad WCAG 2.1 AA
\end{bashcode}

\subsection{Quality gate del proyecto}

El proyecto mantiene un \textit{quality gate} completo que debe superarse antes de
cualquier \textit{merge} a \texttt{develop} o \texttt{main}. Se ejecuta con un único
comando:

\begin{bashcode}{Quality gate completo}
make ci   # PHPStan + Psalm + PHPUnit + cs-check
\end{bashcode}

Las herramientas que componen el gate son:

\begin{description}[nosep]
  \item[PHPStan nivel 5] Análisis estático con \texttt{phpstan-baseline.neon}.
    Detecta tipos incorrectos, llamadas a métodos inexistentes y accesos a propiedades
    indefinidas.
  \item[Psalm] Segunda pasada de análisis estático. Complementa PHPStan detectando
    vulnerabilidades de tipo y problemas de flujo de datos.
  \item[PSR-12] Verificación de estilo de código con PHP-CS-Fixer.
    Corrección automática disponible mediante \texttt{make cs-fix}.
  \item[PHPUnit 13.x] Suite completa de tests unitarios e integración.
    Los tests unitarios se ejecutan en paralelo con \texttt{make test-unit}.
    La cobertura se genera con \texttt{make test-coverage} (HTML + Clover).
\end{description}

\begin{infobox}
  El entorno \gls{ci} ejecuta \texttt{make ci} en cada \textit{pull request}.
  Un quality gate fallido bloquea el merge hasta que todos los errores sean corregidos.
\end{infobox}

\section{Estructura de directorios}
\label{sec:estructura}

\begin{figure}[H]
  \centering
  \begin{forest}
    for tree={
      font=\ttfamily\small,
      grow'=0,
      child anchor=west,
      parent anchor=south,
      anchor=west,
      calign=first,
      edge path={
        \noexpand\path [draw, \forestoption{edge}]
        (!u.south west) +(7.5pt,0) |- node[fill,inner sep=1.25pt] {} (.child anchor)\forestoption{edge label};
      },
      before typesetting nodes={
        if n=1 {insert before={[,phantom]}} {}
      },
      fit=band,
      before computing xy={l=15pt},
    }
    [komorebi/
      [app/
        [Core/ [Router, Container, View, DB, Cache, Queue, Logger]]
        [Http/ [Controllers/ [Admin/ Manager/ Kitchen/ Keeper/ ...]]]
        [Services/ [53 servicios de negocio]]
        [Repositories/ [29 repositorios PDO]]
        [Events/ + Listeners/]
        [routes.php]
      ]
      [bootstrap/ [container.php]]
      [migrations/ [001 -- 019 SQL]]
      [resources/ [views/ [layouts/ admin/ public/ ...]]]
      [tests/ [Unit/ Integration/ e2e/]]
      [public/ [index.php css/ js/ images/]]
      [docs/ [diagrams/ images/ screenshots/ plans/]]
    ]
  \end{forest}
  \caption{Estructura de directorios del proyecto}
  \label{fig:estructura}
\end{figure}

\section{Despliegue}
\label{sec:despliegue}

\subsection{Entorno local con Docker Compose}

\begin{bashcode}{Arrancar el entorno de desarrollo}
make dev          # Levanta app (FrankenPHP) + MySQL + Redis + Mailpit
make db-migrate   # Aplica las 19 migraciones SQL
make db-seed      # Carga datos de prueba
make test         # Quality gate completo
\end{bashcode}

El archivo \texttt{docker-compose.yml} define cuatro servicios:
\begin{itemize}[nosep]
  \item \texttt{app} — FrankenPHP (Caddy + PHP 8.4 integrado), puerto 8080
  \item \texttt{mysql} — MySQL 8.0, puerto 3306
  \item \texttt{redis} — Redis 7, puerto 6379
  \item \texttt{mailpit} — SMTP local + UI, puertos 1025/8025
\end{itemize}

\subsection{Producción en Railway}

El despliegue a producción se realiza automáticamente mediante GitHub Actions al
hacer \textit{push} a \texttt{main}. Railway detecta el \texttt{Dockerfile}
(multi-stage build) y despliega la imagen resultante.

Las variables de entorno sensibles se gestionan como secretos en Railway. El helper
\texttt{SecretLoader::require()} busca la variable de entorno primero y después el
archivo \texttt{/run/secrets/<clave>} (estándar Docker Swarm).

\section{Control de versiones y Git Flow}
\label{sec:gitflow}

\begin{figure}[H]
  \centering
  \begin{tikzpicture}[
    commit/.style={draw,circle,fill=komorebi,minimum size=0.6cm,
                   inner sep=0pt,font=\scriptsize\color{white}},
    branch/.style={font=\small\bfseries,komorebi-dark}
  ]
    \node[branch] at (0,3) {main};
    \node[branch] at (0,2) {develop};
    \node[branch] at (0,1) {feature/reservas};
    \node[branch] at (0,0) {hotfix/csrf};

    \foreach \x in {1,3,5,7,9} \node[commit] at (\x,3) {};
    \foreach \x in {1,2,3,4,5,6,7,8,9} \node[commit] at (\x,2) {};
    \foreach \x in {2,3,4,5,6} \node[commit] at (\x,1) {};
    \foreach \x in {4,5} \node[commit] at (\x,0) {};

    \draw[thick,komorebi] (1,3) -- (9,3);
    \draw[thick,komorebi!60] (1,2) -- (9,2);
    \draw[thick,gray] (2,1) -- (6,1);
    \draw[thick,orange!70!black] (4,0) -- (5,0);

    \draw[->,dashed] (6,1) -- (7,2);  % feature merge
    \draw[->,dashed] (5,0) -- (5,2);  % hotfix merge
    \draw[->,dashed] (5,2) -- (5,3);  % hotfix a main
    \draw[->,dashed] (2,2) -- (2,1);  % branch feature
    \draw[->,dashed] (4,2) -- (4,0);  % branch hotfix
  \end{tikzpicture}
  \caption{Flujo de ramas Git Flow del proyecto}
  \label{fig:gitflow}
\end{figure}

% ============================================================
%  CAPÍTULO 3 — MANUAL DE USUARIO
% ============================================================
\chapter{Manual de Usuario}
\label{chap:manual}

Este capítulo describe las funcionalidades de la plataforma desde la perspectiva del
usuario final, organizadas por tipo de acceso y rol. Cada sección incluye capturas
de pantalla reales del sistema.

\begin{infobox}
  Las capturas de pantalla están numeradas del \texttt{01} al \texttt{62} y se
  encuentran en \texttt{docs/screenshots/}. Se referencian mediante
  \verb|\cref{fig:NN-nombre}| a lo largo del capítulo.
\end{infobox}

\section{Área pública}
\label{sec:manual-publico}

El área pública es accesible sin autenticación y muestra la presentación de la cadena,
los cafés disponibles, la carta del menú, el sistema de reservas, la historia de la
marca y los recursos informativos.

\subsection{Página de inicio}
La página principal (\cref{fig:01-home}) presenta el concepto de la marca mediante
un encabezado animado con el significado de \textit{komorebi}, seguido de las
secciones destacadas: cafés disponibles, menú del momento y acceso rápido a reservas.

\begin{figure}[H]
  \centering
  \fbox{\rule{0pt}{6cm}\rule{12cm}{0pt}}
  \caption{Página de inicio — Komorebi Café}
  \label{fig:01-home}
\end{figure}

\subsection{Catálogo de cafés}
La página \texttt{/cafes} (\cref{fig:02-cafes-lista}) muestra las tarjetas de los
catorce cafés temáticos con imagen, nombre y botón de acceso al detalle.

\begin{figure}[H]
  \centering
  \fbox{\rule{0pt}{6cm}\rule{12cm}{0pt}}
  \caption{Catálogo de cafés temáticos}
  \label{fig:02-cafes-lista}
\end{figure}

\subsection{Detalle de café}
Cada café dispone de su propia página de detalle (\cref{fig:03-cafe-detalle}) con
galería de imágenes, descripción temática, horarios, capacidad y acceso directo al
formulario de reserva.

\begin{figure}[H]
  \centering
  \fbox{\rule{0pt}{6cm}\rule{12cm}{0pt}}
  \caption{Página de detalle de un café temático}
  \label{fig:03-cafe-detalle}
\end{figure}

\subsection{Carta del menú}
La carta digital (\cref{fig:04-menu-catalogo}) presenta todos los productos con
foto, descripción y precio. Incluye filtros por categoría y por alérgeno en tiempo
real mediante Alpine.js.

\begin{figure}[H]
  \centering
  \fbox{\rule{0pt}{6cm}\rule{12cm}{0pt}}
  \caption{Carta digital del menú con filtros por alérgeno}
  \label{fig:04-menu-catalogo}
\end{figure}

\subsection{Quiz de recomendación}
El quiz (\cref{fig:06-quiz-inicio}) guía al visitante a través de una serie de
preguntas sobre preferencias estéticas y recomienda el café temático más afín a su
personalidad.

\begin{figure}[H]
  \centering
  \fbox{\rule{0pt}{6cm}\rule{12cm}{0pt}}
  \caption{Quiz de recomendación de café}
  \label{fig:06-quiz-inicio}
\end{figure}

\subsection{Historia de la marca}
La página \texttt{/historia} (\cref{fig:08-historia}) narra el origen del concepto
\textit{komorebi} y la filosofía que inspira cada uno de los establecimientos.

\begin{figure}[H]
  \centering
  \fbox{\rule{0pt}{6cm}\rule{12cm}{0pt}}
  \caption{Página de historia de la marca}
  \label{fig:08-historia}
\end{figure}

\section{Autenticación y gestión de cuenta}
\label{sec:manual-auth}

\subsection{Formulario de login}
El formulario de acceso (\cref{fig:13-login-form}) solicita email y contraseña.
Incluye protección \gls{csrf} y límite de velocidad de 5 intentos por minuto.
El campo contraseña dispone de botón de alternancia de visibilidad.

\begin{figure}[H]
  \centering
  \fbox{\rule{0pt}{6cm}\rule{12cm}{0pt}}
  \caption{Formulario de inicio de sesión}
  \label{fig:13-login-form}
\end{figure}

\subsection{Registro de nueva cuenta}
El formulario de registro (\cref{fig:14-registro-form}) recoge nombre, email y
contraseña. Tras el envío, el sistema envía un email de verificación.

\begin{figure}[H]
  \centering
  \fbox{\rule{0pt}{6cm}\rule{12cm}{0pt}}
  \caption{Formulario de registro de nueva cuenta}
  \label{fig:14-registro-form}
\end{figure}

\subsection{Recuperación de contraseña}
El flujo de recuperación (\cref{fig:15-forgot-password}) envía un enlace de
restablecimiento con validez de 60 minutos al email registrado.

\begin{figure}[H]
  \centering
  \fbox{\rule{0pt}{6cm}\rule{12cm}{0pt}}
  \caption{Formulario de recuperación de contraseña}
  \label{fig:15-forgot-password}
\end{figure}

\section{Usuario autenticado}
\label{sec:manual-usuario}

Una vez autenticado, el usuario accede a su panel personal con reservas, programa de
fidelización, favoritos y carrito.

\subsection{Perfil de usuario}
La página \texttt{/perfil} (\cref{fig:19-perfil}) permite actualizar nombre, email,
avatar y contraseña. Los cambios de contraseña requieren confirmación de la contraseña
actual.

\begin{figure}[H]
  \centering
  \fbox{\rule{0pt}{6cm}\rule{12cm}{0pt}}
  \caption{Página de perfil de usuario}
  \label{fig:19-perfil}
\end{figure}

\subsection{Mis reservas}
La sección de reservas (\cref{fig:22-mis-reservas}) lista todas las reservas del
usuario con estado (pendiente, confirmada, cancelada), fecha, café y número de
comensales. Permite cancelar reservas con más de 24 horas de antelación.

\begin{figure}[H]
  \centering
  \fbox{\rule{0pt}{6cm}\rule{12cm}{0pt}}
  \caption{Listado de reservas del usuario}
  \label{fig:22-mis-reservas}
\end{figure}

\subsection{Nueva reserva}
El asistente de reserva en dos pasos (\cref{fig:23-reservar-step1}) muestra
primero la selección de café, fecha y franja horaria disponible. En el segundo
paso (\cref{fig:24-reservar-step2}) se confirma el número de comensales y las
preferencias dietéticas.

\begin{figure}[H]
  \centering
  \fbox{\rule{0pt}{6cm}\rule{12cm}{0pt}}
  \caption{Formulario de reserva — paso 1}
  \label{fig:23-reservar-step1}
\end{figure}

\begin{figure}[H]
  \centering
  \fbox{\rule{0pt}{6cm}\rule{12cm}{0pt}}
  \caption{Formulario de reserva — paso 2}
  \label{fig:24-reservar-step2}
\end{figure}

\subsection{Tarjeta de fidelización}
La tarjeta digital (\cref{fig:26-loyalty-card}) muestra los puntos acumulados,
el nivel (\textit{bronze}, \textit{silver}, \textit{gold} o \textit{platinum}), el historial de
transacciones y los premios canjeables del catálogo.

\begin{figure}[H]
  \centering
  \fbox{\rule{0pt}{6cm}\rule{12cm}{0pt}}
  \caption{Tarjeta de fidelización digital}
  \label{fig:26-loyalty-card}
\end{figure}

\section{Panel de administración}
\label{sec:manual-admin}

El backoffice de administración (\texttt{/admin}) está reservado al rol
\texttt{admin} y proporciona control total sobre todas las entidades del sistema.

\subsection{Dashboard de administración}
El dashboard (\cref{fig:30-admin-dashboard}) muestra KPIs en tiempo real:
reservas del día, ingresos semanales, nuevos usuarios, alertas de animales y
estado del sistema.

\begin{figure}[H]
  \centering
  \fbox{\rule{0pt}{6cm}\rule{12cm}{0pt}}
  \caption{Dashboard de administración}
  \label{fig:30-admin-dashboard}
\end{figure}

\subsection{Gestión de usuarios}
La tabla de usuarios (\cref{fig:31-admin-usuarios}) permite buscar, filtrar por
rol, activar/desactivar cuentas y acceder al detalle de cada usuario.

\begin{figure}[H]
  \centering
  \fbox{\rule{0pt}{6cm}\rule{12cm}{0pt}}
  \caption{Panel de gestión de usuarios}
  \label{fig:31-admin-usuarios}
\end{figure}

\subsection{Gestión de cafés}
La administración de cafés (\cref{fig:32-admin-cafes}) permite crear nuevos
establecimientos, editar sus datos (capacidad, horarios, temática) y activar o
desactivar su presencia pública.

\begin{figure}[H]
  \centering
  \fbox{\rule{0pt}{6cm}\rule{12cm}{0pt}}
  \caption{Panel de gestión de cafés}
  \label{fig:32-admin-cafes}
\end{figure}

\subsection{Gestión del menú}
La carta digital (\cref{fig:33-admin-menu}) permite administrar categorías y
productos, definir precios, alérgenos y disponibilidad por café.

\begin{figure}[H]
  \centering
  \fbox{\rule{0pt}{6cm}\rule{12cm}{0pt}}
  \caption{Panel de gestión del menú}
  \label{fig:33-admin-menu}
\end{figure}

\subsection{Gestión de reservas}
El panel de reservas (\cref{fig:34-admin-reservas}) ofrece una vista tabular de
todas las reservas con filtros por café, fecha, estado y usuario. Permite confirmar
o cancelar reservas individualmente.

\begin{figure}[H]
  \centering
  \fbox{\rule{0pt}{6cm}\rule{12cm}{0pt}}
  \caption{Panel de gestión de reservas}
  \label{fig:34-admin-reservas}
\end{figure}

\subsection{Gestión de animales}
La sección de animales (\cref{fig:35-admin-animales}) lista todos los residentes
de los cafés con su especie, estado de salud, café de origen y acceso al historial
de incidentes y revisiones.

\begin{figure}[H]
  \centering
  \fbox{\rule{0pt}{6cm}\rule{12cm}{0pt}}
  \caption{Panel de gestión de animales}
  \label{fig:35-admin-animales}
\end{figure}

\subsection{Logs de auditoría}
El visor de logs (\cref{fig:41-admin-logs-auditoria}) registra todas las acciones
administrativas con usuario, timestamp, tipo de acción y entidad afectada.
Permite exportar en CSV.

\begin{figure}[H]
  \centering
  \fbox{\rule{0pt}{6cm}\rule{12cm}{0pt}}
  \caption{Visor de logs de auditoría}
  \label{fig:41-admin-logs-auditoria}
\end{figure}

\subsection{Configuración del sistema}
La pantalla de configuración (\cref{fig:40-admin-settings}) agrupa los parámetros
del sistema por grupos: general, email, reservas, fidelización y mantenimiento.

\begin{figure}[H]
  \centering
  \fbox{\rule{0pt}{6cm}\rule{12cm}{0pt}}
  \caption{Pantalla de configuración del sistema}
  \label{fig:40-admin-settings}
\end{figure}

\section{Panel de manager}
\label{sec:manual-manager}

El manager gestiona las operaciones diarias de su café asignado: productos,
staff, reservas e informes.

\subsection{Dashboard de manager}
El dashboard (\cref{fig:49-manager-dashboard}) muestra el resumen del día:
reservas confirmadas, ocupación actual, ingresos y alertas de personal.

\begin{figure}[H]
  \centering
  \fbox{\rule{0pt}{6cm}\rule{12cm}{0pt}}
  \caption{Dashboard del manager}
  \label{fig:49-manager-dashboard}
\end{figure}

\subsection{Gestión de productos}
La carta del manager (\cref{fig:50-manager-productos}) permite activar/desactivar
productos específicos de su café sin necesidad de acudir al administrador.

\begin{figure}[H]
  \centering
  \fbox{\rule{0pt}{6cm}\rule{12cm}{0pt}}
  \caption{Gestión de productos del manager}
  \label{fig:50-manager-productos}
\end{figure}

\section{Recepción y Cocina}
\label{sec:manual-ops}

\subsection{Panel de recepción}
El panel de recepción (\cref{fig:54-reception-dashboard}) muestra en tiempo real
las reservas del día con sus estados: esperando, sentada, en consumo, pagada.
Permite hacer check-in con un clic.

\begin{figure}[H]
  \centering
  \fbox{\rule{0pt}{6cm}\rule{12cm}{0pt}}
  \caption{Panel de recepción con reservas del día}
  \label{fig:54-reception-dashboard}
\end{figure}

\subsection{Kitchen Display System (KDS)}
El \gls{kds} (\cref{fig:56-kitchen-orders}) muestra las órdenes activas en tiempo
real mediante Server-Sent Events (Mercure). Cada tarjeta incluye el número de mesa,
los ítems y el tiempo transcurrido desde la entrada del pedido.

\begin{figure}[H]
  \centering
  \fbox{\rule{0pt}{6cm}\rule{12cm}{0pt}}
  \caption{Kitchen Display System — órdenes activas}
  \label{fig:56-kitchen-orders}
\end{figure}

\section{Panel del Keeper}
\label{sec:manual-keeper}

El rol \texttt{keeper} gestiona el bienestar de los animales residentes de los cafés
temáticos.

\subsection{Dashboard del Keeper}
El dashboard (\cref{fig:58-keeper-dashboard}) muestra el estado general de todos
los animales activos: revisiones pendientes, incidentes abiertos y alertas de salud.

\begin{figure}[H]
  \centering
  \fbox{\rule{0pt}{6cm}\rule{12cm}{0pt}}
  \caption{Dashboard del Keeper con estado de animales}
  \label{fig:58-keeper-dashboard}
\end{figure}

\subsection{Ficha de animal}
La ficha individual (\cref{fig:59-keeper-animals}) incluye datos biográficos,
historial de revisiones, incidentes pasados y foto del animal.

\begin{figure}[H]
  \centering
  \fbox{\rule{0pt}{6cm}\rule{12cm}{0pt}}
  \caption{Catálogo de animales del café}
  \label{fig:59-keeper-animals}
\end{figure}

\subsection{Revisiones de salud}
El formulario de revisión (\cref{fig:61-keeper-health-check-crear}) registra el
estado general, peso, temperatura, observaciones y la fecha de la próxima revisión.

\begin{figure}[H]
  \centering
  \fbox{\rule{0pt}{6cm}\rule{12cm}{0pt}}
  \caption{Formulario de nueva revisión de salud}
  \label{fig:61-keeper-health-check-crear}
\end{figure}

\subsection{Registro de incidentes}
El historial de incidentes (\cref{fig:62-keeper-incidents}) categoriza cada evento
por severidad (\texttt{low}, \texttt{medium}, \texttt{high}, \texttt{critical}),
su descripción, la resolución aplicada y la fecha de cierre.

\begin{figure}[H]
  \centering
  \fbox{\rule{0pt}{6cm}\rule{12cm}{0pt}}
  \caption{Registro de incidentes de animales}
  \label{fig:62-keeper-incidents}
\end{figure}

% ============================================================
%  CAPÍTULO 4 — CONCLUSIONES Y FUTURAS MEJORAS
% ============================================================
\chapter{Conclusiones y Futuras Mejoras}
\label{chap:conclusiones}

\section{Objetivos alcanzados}
\label{sec:objetivos-alcanzados}

Al inicio del proyecto se plantearon los siguientes objetivos principales, todos
ellos cumplidos en la versión actual:

\begin{enumerate}
  \item \textbf{Plataforma funcional de gestión integral}: el sistema cubre el ciclo
    completo de operaciones de la cadena, desde la reserva pública hasta el cierre
    de caja en recepción.
  \item \textbf{Framework MVC propio sobre PHP~8.4}: se ha demostrado la comprensión
    profunda de los mecanismos subyacentes al construir el router, el contenedor IoC,
    el motor de vistas y el sistema de middleware desde cero.
  \item \textbf{Arquitectura orientada a estándares}: adherencia completa a PSR-7,
    PSR-14 y PSR-15; Result pattern en todos los servicios; DTOs \texttt{final readonly};
    inyección por interfaz en todos los controladores.
  \item \textbf{Calidad verificada}: PHPStan nivel~5, Psalm, suite PHPUnit con tests
    unitarios y de integración, tests E2E con Playwright, análisis de cobertura.
  \item \textbf{Despliegue en producción}: imagen Docker multi-stage desplegada en
    Railway con GitHub Actions como pipeline \gls{ci}.
  \item \textbf{Identidad visual coherente}: el diseño de la interfaz refleja la
    estética \textit{komorebi}: paleta natural, tipografía serena y microanimaciones
    evocadoras.
\end{enumerate}

\section{Aprendizajes técnicos}
\label{sec:aprendizajes}

El desarrollo de Komorebi Café ha permitido profundizar en áreas que habitualmente
quedan ocultas detrás de los frameworks de alto nivel:

\begin{description}[style=nextline]
  \item[PSR-7 y el ciclo petición-respuesta]
    Trabajar directamente con \texttt{ServerRequestInterface} y \texttt{ResponseInterface}
    obliga a entender qué hace realmente un framework cuando procesa una petición HTTP.
  \item[El Result pattern como alternativa a las excepciones]
    Diseñar todos los servicios para que devuelvan \texttt{Result::ok()/fail()} en
    lugar de lanzar excepciones para fallos esperados ha simplificado el flujo de
    control y hecho el código más predecible.
  \item[Redis como infraestructura transversal]
    Usar Redis no solo para caché sino también para colas de jobs, sesiones,
    rate-limiting e idempotencia ha demostrado la versatilidad de esta herramienta.
  \item[12-Factor en la práctica]
    Externalizar toda la configuración en variables de entorno, escribir logs a
    stdout y construir la imagen Docker de forma que sea idéntica en dev y producción
    simplifica enormemente el ciclo de despliegue.
  \item[Accesibilidad como requisito, no como mejora]
    Integrar los tests de accesibilidad WCAG~2.1 AA en el pipeline desde el principio
    evita deuda técnica acumulada y garantiza una experiencia inclusiva.
\end{description}

\section{Futuras mejoras propuestas}
\label{sec:futuras}

\begin{enumerate}
  \item \textbf{Aplicación móvil nativa}: iOS + Android consumiendo la \gls{api} REST
    ya disponible, con notificaciones push mediante Firebase Cloud Messaging.
  \item \textbf{Motor de recomendación}: sistema de recomendación de cafés y productos
    basado en el historial de reservas y las preferencias del programa de fidelización.
  \item \textbf{Integración con pasarela de pago}: Stripe o Redsys para cobro online
    en el momento de la reserva.
  \item \textbf{Panel de analytics avanzado}: integración con Grafana y Prometheus
    para monitorización de métricas de negocio en tiempo real.
  \item \textbf{Internacionalización (i18n)}: soporte multiidioma (español, inglés,
    japonés) para el área pública y la carta digital.
  \item \textbf{Tests de mutación}: implementar Infection PHP para medir la efectividad
    real de la suite de tests unitarios.
  \item \textbf{GraphQL para la API pública}: complementar la API REST con un endpoint
    GraphQL que permita queries más flexibles desde aplicaciones móviles.
  \item \textbf{Modo offline en el KDS}: Service Worker para que el \gls{kds} de cocina
    siga operativo durante cortes de conectividad.
\end{enumerate}

% ============================================================
%  CAPÍTULO 5 — REFERENCIAS BIBLIOGRÁFICAS
% ============================================================
\chapter{Referencias Bibliográficas}
\label{chap:refs}

\printbibliography[heading=none]

% ============================================================
%  BACKMATTER — APÉNDICES
% ============================================================
% \backmatter  % solo disponible en clase book — no existe en report

\begin{appendices}

% ─── APÉNDICE A ──────────────────────────────────────────────
\chapter{Catálogo completo de endpoints API REST}
\label{ap:api}

La \cref{tab:api-completo} recoge todos los endpoints de la API REST pública y
privada de Komorebi Café. Los endpoints marcados con \dag\ requieren autenticación
Bearer token.

\begin{landscape}
\begin{longtable}{@{}lllp{5cm}@{}}
  \caption{Catálogo completo de endpoints API REST (/api/v1/)}
  \label{tab:api-completo}\\
  \toprule
  \textbf{Método} & \textbf{Ruta} & \textbf{Auth} & \textbf{Descripción}\\
  \midrule
  \endfirsthead
  \toprule
  \textbf{Método} & \textbf{Ruta} & \textbf{Auth} & \textbf{Descripción}\\
  \midrule
  \endhead
  \bottomrule
  \endfoot
  % ── Menú público ──
  GET  & \texttt{/menu/alergenos}            & —    & Lista de alérgenos\\
  GET  & \texttt{/menu/productos}            & —    & Catálogo de productos\\
  GET  & \texttt{/menu/products/\{id\}}      & —    & Detalle de producto\\
  % ── Cafés público ──
  GET  & \texttt{/cafes}                     & —    & Lista de cafés activos\\
  GET  & \texttt{/cafes/\{slug\}}            & —    & Detalle de café\\
  % ── Festivos y disponibilidad ──
  GET  & \texttt{/holidays}                  & —    & Festivos del año\\
  GET  & \texttt{/holidays/\{date\}}         & —    & Verifica si una fecha es festivo\\
  GET  & \texttt{/time-slots/available}      & —    & Franjas disponibles por café y fecha\\
  GET  & \texttt{/time-slots/stats}          & —    & Estadísticas de ocupación\\
  % ── Pases ──
  GET  & \texttt{/passes}                    & —    & Tipos de pase disponibles\\
  % ── Cookies / preferencias ──
  PATCH& \texttt{/cookies}                   & —    & Guardar consentimiento de cookies\\
  GET  & \texttt{/cookies/filters}           & —    & Obtener filtros de dieta\\
  PUT  & \texttt{/cookies/filters}           & —    & Guardar filtros de dieta\\
  DELETE&\texttt{/cookies/filters}           & —    & Eliminar filtros\\
  PUT  & \texttt{/cookies/dietary}           & —    & Guardar preferencias dietéticas\\
  POST & \texttt{/cookies/recently-viewed}   & —    & Añadir producto a recientes\\
  GET  & \texttt{/cookies/recently-viewed/data}& —  & Obtener datos de recientes\\
  DELETE&\texttt{/cookies/recently-viewed}   & —    & Limpiar recientes\\
  % ── Newsletter ──
  POST & \texttt{/newsletter/subscriptions}  & —    & Suscribirse al boletín\\
  % ── Waitlist pública ──
  GET  & \texttt{/waitlists/\{token\}}       & —    & Posición en lista de espera\\
  POST & \texttt{/waitlists/\{token\}/confirmations}& —& Confirmar posición\\
  POST & \texttt{/waitlists}                 & \dag & Unirse a lista de espera\\
  % ── Tokens ──
  GET  & \texttt{/tokens}                    & \dag & Listar tokens propios\\
  POST & \texttt{/tokens}                    & —    & Crear token Bearer\\
  DELETE&\texttt{/tokens/\{id\}}             & \dag & Revocar token\\
  % ── Favoritos ──
  GET  & \texttt{/favorites}                 & \dag & Listar favoritos\\
  PUT  & \texttt{/favorites/\{id\}}          & \dag & Añadir favorito\\
  DELETE&\texttt{/favorites/\{id\}}          & \dag & Eliminar favorito\\
  % ── Carrito ──
  GET  & \texttt{/cart}                      & \dag & Obtener carrito\\
  GET  & \texttt{/cart/guest}                & —    & Carrito de invitado\\
  POST & \texttt{/cart/items}                & \dag & Añadir ítem\\
  PATCH& \texttt{/cart/items/\{id\}}         & \dag & Actualizar cantidad\\
  DELETE&\texttt{/cart/items/\{id\}}         & \dag & Eliminar ítem\\
  DELETE&\texttt{/cart/items}                & \dag & Vaciar carrito\\
  % ── Reservas ──
  GET  & \texttt{/reservations/available}    & —    & Huecos disponibles\\
  POST & \texttt{/reservations}              & \dag & Crear reserva (idempotente)\\
  GET  & \texttt{/reservations/\{id\}}       & \dag & Detalle de reserva\\
  POST & \texttt{/reservations/\{id\}/cancel}& \dag & Cancelar reserva\\
  GET  & \texttt{/user/reservations}         & \dag & Mis reservas\\
  % ── Fidelización ──
  GET  & \texttt{/loyalty/validate/\{code\}} & \dag & Validar código de canje\\
  POST & \texttt{/loyalty/use}               & \dag & Usar puntos\\
  POST & \texttt{/loyalty/redeem}            & \dag & Canjear recompensa\\
  % ── Reseñas ──
  POST & \texttt{/reviews}                   & \dag & Crear reseña\\
  PUT  & \texttt{/reviews/\{id\}}            & \dag & Actualizar reseña\\
  DELETE&\texttt{/reviews/\{id\}}            & \dag & Eliminar reseña\\
  % ── Usuario ──
  GET  & \texttt{/user/profile}              & \dag & Perfil del usuario\\
  GET  & \texttt{/user/stats}                & \dag & Estadísticas del usuario\\
  GET  & \texttt{/user/reviews}              & \dag & Reseñas del usuario\\
  GET  & \texttt{/user/avatar-options}       & \dag & Opciones de avatar\\
  % ── Manager ──
  GET  & \texttt{/manager/stats}             & \dag & KPIs del manager\\
  GET  & \texttt{/manager/revenue/weekly}    & \dag & Ingresos semanales\\
  GET  & \texttt{/manager/animals/top}       & \dag & Animales más populares\\
  GET  & \texttt{/manager/reservations/status}& \dag & Estado reservas del café\\
  % ── Admin — Usuarios ──
  POST & \texttt{/admin/users}               & \dag & Crear usuario\\
  PUT  & \texttt{/admin/users/\{id\}}        & \dag & Actualizar usuario\\
  DELETE&\texttt{/admin/users/\{id\}}        & \dag & Eliminar usuario\\
  PATCH& \texttt{/admin/users/\{id\}/status} & \dag & Toggle activo/inactivo\\
  % ── Admin — Cafés ──
  POST & \texttt{/admin/cafes}               & \dag & Crear café\\
  PUT  & \texttt{/admin/cafes/\{id\}}        & \dag & Actualizar café\\
  DELETE&\texttt{/admin/cafes/\{id\}}        & \dag & Eliminar café\\
  PATCH& \texttt{/admin/cafes/\{id\}/status} & \dag & Toggle estado café\\
  % ── Admin — Menú ──
  POST & \texttt{/admin/menu}                & \dag & Crear producto\\
  PUT  & \texttt{/admin/menu/\{id\}}         & \dag & Actualizar producto\\
  DELETE&\texttt{/admin/menu/\{id\}}         & \dag & Eliminar producto\\
  PATCH& \texttt{/admin/menu/\{id\}/toggle}  & \dag & Toggle disponibilidad\\
  % ── Admin — Reservas ──
  POST & \texttt{/admin/reservations/\{id\}/confirm}& \dag & Confirmar reserva\\
  POST & \texttt{/admin/reservations/\{id\}/cancel} & \dag & Cancelar reserva (admin)\\
  % ── Admin — Reseñas ──
  POST & \texttt{/admin/reviews/\{id\}/approve}& \dag & Aprobar reseña\\
  POST & \texttt{/admin/reviews/\{id\}/reject} & \dag & Rechazar reseña\\
  DELETE&\texttt{/admin/reviews/\{id\}}       & \dag & Eliminar reseña (admin)\\
  % ── Admin — Newsletter ──
  DELETE&\texttt{/admin/newsletter/subscribers/\{email\}}& \dag & Dar de baja\\
  GET  & \texttt{/admin/newsletter/export}   & \dag & Exportar suscriptores CSV\\
  % ── Admin — Fidelización ──
  GET  & \texttt{/admin/loyalty/stats}       & \dag & Estadísticas de fidelización\\
  GET  & \texttt{/admin/loyalty/cards}       & \dag & Listado de tarjetas\\
  GET  & \texttt{/admin/loyalty/catalog}     & \dag & Catálogo de recompensas\\
  PATCH& \texttt{/admin/loyalty/catalog/\{id\}/toggle}& \dag & Toggle recompensa\\
  GET  & \texttt{/admin/loyalty/redemptions} & \dag & Historial de canjes\\
  % ── Admin — Logs ──
  GET  & \texttt{/admin/logs/audit}          & \dag & Logs de auditoría\\
  GET  & \texttt{/admin/logs/audit/export}   & \dag & Exportar audit log CSV\\
  GET  & \texttt{/admin/logs/auth}           & \dag & Logs de autenticación\\
  GET  & \texttt{/admin/logs/auth/suspicious-count}& \dag & Contador de accesos sospechosos\\
  GET  & \texttt{/admin/logs/auth/export}    & \dag & Exportar auth log CSV\\
  % ── Admin — Sistema ──
  GET  & \texttt{/admin/settings}            & \dag & Obtener configuración\\
  PUT  & \texttt{/admin/settings/\{group\}}  & \dag & Actualizar grupo de config.\\
  POST & \texttt{/admin/settings/test-email} & \dag & Enviar email de prueba\\
  POST & \texttt{/admin/cache/clear}         & \dag & Vaciar caché Redis\\
  POST & \texttt{/admin/security/block-ip}   & \dag & Bloquear IP\\
  % ── Operaciones: Recepción ──
  GET  & \texttt{/ops/reception/reservations}          & \dag & Reservas del día\\
  POST & \texttt{/ops/reception/reservations/\{id\}/checkin}   & \dag & Check-in\\
  POST & \texttt{/ops/reception/reservations/\{id\}/checkout}  & \dag & Check-out\\
  POST & \texttt{/ops/reception/reservations/\{id\}/items}     & \dag & Añadir ítems consumidos\\
  POST & \texttt{/ops/reception/reservations/\{id\}/payments}  & \dag & Registrar pago\\
  % ── Operaciones: Cocina ──
  GET  & \texttt{/ops/kitchen/orders}                & \dag & Órdenes activas\\
  POST & \texttt{/ops/kitchen/orders/\{id\}/complete}& \dag & Marcar orden como completada\\
  % ── Supervisor ──
  POST & \texttt{/supervisor/assignments}    & \dag & Crear asignación\\
  GET  & \texttt{/supervisor/assignments}    & \dag & Listar asignaciones\\
\end{longtable}
\end{landscape}

% ─── APÉNDICE B ──────────────────────────────────────────────
\chapter{Esquema de base de datos}
\label{ap:bd}

La base de datos de Komorebi Café está definida en 19 archivos de migración SQL
ubicados en \texttt{migrations/}. La \cref{tab:migraciones} describe cada migración.

\begin{longtable}{@{}llp{7cm}@{}}
  \caption{Historial de migraciones SQL}
  \label{tab:migraciones}\\
  \toprule
  \textbf{Nº} & \textbf{Archivo} & \textbf{Contenido}\\
  \midrule
  \endfirsthead
  \toprule
  \textbf{Nº} & \textbf{Archivo} & \textbf{Contenido}\\
  \midrule
  \endhead
  \bottomrule
  \endfoot
  001 & \texttt{001\_infrastructure.sql}       & Infraestructura base: tablas de configuración\\
  002 & \texttt{002\_users\_rbac.sql}           & Usuarios, roles, permisos RBAC\\
  003 & \texttt{003\_reviews.sql}               & Reseñas de clientes\\
  004 & \texttt{004\_reservations.sql}          & Reservas, ítems y time slots\\
  005 & \texttt{005\_email\_auth.sql}           & Tokens de email y autenticación\\
  006 & \texttt{006\_telegram\_bot.sql}         & Configuración del bot de Telegram\\
  007 & \texttt{007\_external\_cache.sql}       & Caché externo de API (festivos, clima)\\
  008 & \texttt{008\_animals.sql}               & Animales, incidentes y galería\\
  009 & \texttt{009\_system\_settings.sql}      & Configuración del sistema por grupos\\
  010 & \texttt{010\_newsletter.sql}            & Suscriptores del boletín\\
  011 & \texttt{011\_time\_slots\_waitlist.sql} & Franjas horarias y lista de espera\\
  012 & \texttt{012\_waitlist.sql}              & Estructura definitiva de waitlist\\
  012b& \texttt{012b\_reservation\_triggers.sql}& Triggers SQL de reservas\\
  013 & \texttt{013\_loyalty\_system.sql}       & Tarjetas, transacciones y catálogo de fidelización\\
  014 & \texttt{014\_staff\_shifts.sql}         & Turnos de personal\\
  015 & \texttt{015\_animal\_health\_checks.sql}& Revisiones de salud de animales\\
  016 & \texttt{016\_supervisor\_assignments.sql}& Asignaciones de supervisor\\
  018 & \texttt{018\_api\_tokens.sql}           & Tokens Bearer para API REST\\
  —   & \texttt{026--030\_*.sql}                & Correcciones y ajustes adicionales\\
\end{longtable}

% ─── APÉNDICE C ──────────────────────────────────────────────
\chapter{Variables de entorno}
\label{ap:env}

La configuración de la aplicación sigue el principio 12-Factor III: toda la
configuración se externaliza en variables de entorno, nunca en el código fuente.
La \cref{tab:envvars} describe las variables principales.

\begin{longtable}{@{}lp{4cm}p{5cm}@{}}
  \caption{Variables de entorno del sistema}
  \label{tab:envvars}\\
  \toprule
  \textbf{Variable} & \textbf{Ejemplo} & \textbf{Descripción}\\
  \midrule
  \endfirsthead
  \toprule
  \textbf{Variable} & \textbf{Ejemplo} & \textbf{Descripción}\\
  \midrule
  \endhead
  \bottomrule
  \endfoot
  \texttt{APP\_ENV}          & \texttt{production}      & Entorno de ejecución\\
  \texttt{APP\_DEBUG}        & \texttt{false}            & Modo debug (nunca true en prod.)\\
  \texttt{APP\_URL}          & \texttt{https://...}      & URL base pública\\
  \texttt{DB\_HOST}          & \texttt{mysql}            & Host de la base de datos\\
  \texttt{DB\_DATABASE}      & \texttt{komorebi}         & Nombre de la base de datos\\
  \texttt{DB\_USERNAME}      & \texttt{komorebi\_user}   & Usuario MySQL\\
  \texttt{DB\_PASSWORD}      & (secreto)                 & Contraseña MySQL\\
  \texttt{REDIS\_HOST}       & \texttt{redis}            & Host Redis\\
  \texttt{REDIS\_PORT}       & \texttt{6379}             & Puerto Redis\\
  \texttt{SESSION\_NAME}     & \texttt{komorebi\_sess}   & Nombre de la cookie de sesión\\
  \texttt{SESSION\_LIFETIME} & \texttt{7200}             & Duración en segundos\\
  \texttt{MAIL\_HOST}        & \texttt{mailpit}          & Servidor SMTP\\
  \texttt{MAIL\_PORT}        & \texttt{1025}             & Puerto SMTP\\
  \texttt{MAIL\_FROM}        & \texttt{no-reply@...}     & Remitente de emails\\
  \texttt{FEATURE\_OPS}      & \texttt{1}                & Activa módulo Ops (recepción/cocina)\\
  \texttt{FEATURE\_BACKOFFICE}& \texttt{1}               & Activa módulo backoffice admin\\
  \texttt{FEATURE\_KEEPER}   & \texttt{1}                & Activa módulo Keeper\\
  \texttt{TELEGRAM\_BOT\_TOKEN}& (secreto)               & Token del bot de Telegram\\
  \texttt{CLOUDINARY\_URL}   & \texttt{cloudinary://...} & URL de Cloudinary (imágenes)\\
\end{longtable}

% ─── APÉNDICE D — GLOSARIO ───────────────────────────────────
\chapter{Glosario de términos y acrónimos}
\label{ap:glosario}

\printnoidxglossary[type=abbreviations,title=Acrónimos]
\printnoidxglossary[type=main,title=Términos]

\end{appendices}

\end{document}
